% Status: überarbeitet als FHNW-Literatur-Review (Begriffe, Theorie, Best
% Practice, Vorbefunde, Forschungslücke). Strukturvorlage angepasst an die
% eigene Arbeit (docs/project/FHNW_ACADEMICGUIDE_RULES.md, Abschnitt 4 + 10).
% Finale Zitation nach Quellenreview (docs/project/SOURCE_REVIEW_LOG.md).
\chapter{Theoretischer Rahmen und Forschungsstand}
\label{ch:theorie}

Dieses Kapitel ordnet die für die Arbeit relevanten Begriffe, Theorien,
Messverfahren und Vorbefunde ein. Es folgt der Logik eines Literatur-Reviews.
Zuerst werden die zentralen Begriffe geklärt, danach die theoretischen
Grundlagen und die etablierten Verfahren zur Messung von Prognosequalität und
Marktreaktion dargestellt, anschliessend die bisherige Forschung zu Polymarket
und zu informiertem Handel diskutiert. Aus dieser Zusammenschau ergibt sich die
Forschungslücke, welche die drei Hypothesen motiviert. Ausgewählt werden nur
jene Konzepte und Quellen, die für die Forschungsfrage unmittelbar relevant sind.

\section{Begriffsklärung}
Ein \emph{Prognosemarkt} ist ein Markt, auf dem Kontrakte auf den Ausgang
künftiger Ereignisse gehandelt werden und dessen Preis sich als kollektive
Wahrscheinlichkeitsschätzung lesen lässt \citep{zotero_poly_005}. Ein binärer
Kontrakt zahlt bei Eintritt des Ereignisses einen festen Betrag und sonst nichts,
sodass der Gleichgewichtspreis zwischen null und eins der vom Markt
eingepreisten Eintrittswahrscheinlichkeit entspricht. \emph{Informationelle
Effizienz} bezeichnet den Grad, in dem dieser Preis die öffentlich verfügbare
Information widerspiegelt \citep{lit_emh_001}. Die Arbeit verwendet den Begriff
bewusst graduell und scope-spezifisch, nicht als absolute Eigenschaft.

\emph{Kalibrierung} meint die Uebereinstimmung von prognostizierter
Wahrscheinlichkeit und beobachteter relativer Häufigkeit: Ereignisse, denen ein
Markt 70 Prozent zuweist, sollten langfristig in etwa 70 Prozent der Fälle
eintreten. \emph{Polymarket} ist ein dezentraler, auf einer Blockchain
abgewickelter Prognosemarkt, bei dem jede Transaktion öffentlich einsehbar ist.
Eine \emph{Wallet} ist die öffentliche Adresse, über die eine Teilnehmerin oder
ein Teilnehmer handelt. Ein \emph{Tier} ist in dieser Arbeit eine
dataset-relative Grössenklasse von Wallets, abgeleitet aus den beobachteten
kumulierten Betragsperzentilen und nicht aus einer willkürlichen
Dollar-Schwelle. Diese Begriffe bilden die gemeinsame Grundlage der drei
Hypothesen.

\section{Markteffizienz und Informationsaggregation}
Der theoretische Ausgangspunkt ist die Markteffizienzhypothese
\citep{lit_emh_001}, nach der Preise verfügbare Information widerspiegeln. Ueblich
ist die Unterscheidung in eine schwache, eine halbstrenge und eine strenge Form,
je nachdem, ob historische Preise, alle öffentlichen Informationen oder auch
private Informationen bereits eingepreist sind. Auf Prognosemärkte übertragen
heisst das: Wenn genügend Teilnehmende ihre Erwartungen laufend an neue Evidenz
anpassen, sollte der Marktpreis als Wahrscheinlichkeit die verfügbare
Information bündeln.

Diesen Aggregationsmechanismus beschreibt \citet{zotero_poly_005} konzeptionell
und verweist auf Vergleiche, in denen spekulative Märkte Experteneinschätzungen
und Umfragen erreichen oder schlagen. Die Idee der Informationsaggregation
(häufig als Weisheit der Vielen bezeichnet) motiviert überhaupt erst, einen
Marktpreis als Forecast zu behandeln. Sie ist zugleich eine Annahme, die
empirisch zu prüfen und nicht vorauszusetzen ist.

\section{Prognosemärkte im Vergleich zu traditionellen Quellen}
Ob Märkte Umfragen tatsächlich überlegen sind, ist empirisch umstritten.
\citet{zotero_poly_002} finden für die US-Wahl 2024 einen Vorteil von Polymarket
gegenüber Umfragen, besonders in den umkämpften Swing States, und ordnen ihn der
Wisdom-of-Crowds-Idee zu. Zugleich betonen sie ausdrücklich, dass die Ergebnisse
weiterer Untersuchung und einer Prüfung auf Uebertragbarkeit bedürfen. Dem steht
eine kritische Sicht gegenüber: \citet{zotero_poly_009} argumentiert, dass
Prognosemärkte keine "Truth Machine" sind, sondern reflexiv sein können ---
Vertrauen in den Marktpreis kann die laufende Informationsaktualisierung
schwächen und so die Genauigkeit untergraben.

Auch die Mikrostruktur solcher Märkte ist relevant. \citet{zotero_poly_006}
analysieren am Beispiel von Kalshi das Zusammenspiel von liquiditätsgebenden und
liquiditätsnehmenden Teilnehmenden und zeigen, dass die Preisqualität von der
Marktstruktur und den Anreizen abhängt. Diese Spannung zwischen belegtem Vorteil
und berechtigter Vorsicht begründet, warum die vorliegende Arbeit nur begrenzte,
scope-spezifische Aussagen anstrebt.

\section{Messung von Prognosequalität}
Zur Messung der Forecast-Qualität dient der Brier-Score \citep{lit_brier_001},
definiert als mittlere quadratische Abweichung zwischen Wahrscheinlichkeitsprognose
und tatsächlichem Ausgang. Ein niedrigerer Brier-Verlust bedeutet eine bessere
Prognose, und der Score belohnt sowohl Treffsicherheit als auch gute Kalibrierung.
Der Vergleich zweier Verlustreihen erfolgt mit dem Diebold-Mariano-Test
\citep{lit_dm_001}, der prüft, ob zwei Prognosen im Mittel gleich genau sind, ohne
einen bestimmten Entstehungsmechanismus zu unterstellen. Beide Verfahren gelten
als etablierte, deterministisch berechenbare Masse der Forecast-Bewertung und
bilden die methodische Grundlage von H1.

\section{Ereignisstudien}
Die Reaktion eines Marktes auf neue Information lässt sich mit Ereignisstudien
untersuchen \citep{lit_eventstudy_001}. Das Verfahren definiert ein Schätzfenster
für das normale Verhalten und ein Ereignisfenster, in dem abnormale und
kumulierte abnormale Bewegungen gemessen werden. Ursprünglich für Aktienrenditen
entwickelt, lässt sich das Design auf Tagesebene und auf Wahrscheinlichkeits-
zeitreihen übertragen. H2 nutzt dieses Design mit vorab fixierten, kuratierten
öffentlichen Ereignissen, um die Reaktionsrichtung von Polymarket nachvollziehbar
zu prüfen.

\section{On-Chain-Aktivität und informierter Handel}
Die Blockchain-Abwicklung macht Wallet-Aktivität sichtbar und eröffnet die Frage,
ob bestimmte Adressen Preisbewegungen zeitlich vorauslaufen. Methodisch nutzt H3
Lead-Lag-Korrelationen und Granger-Tests \citep{lit_granger_001}, deren Ergebnisse
ausdrücklich als prädiktive Timing-Diagnostik und nicht als strukturelle
Kausalität zu lesen sind. Theoretisch rahmt das Modell von \citet{lit_kyle_001}
informierten Handel: Eine informierte Partei verbirgt ihre Order im normalen
Handelsfluss, was sich in abnormaler Handelsgrösse niederschlagen kann. Empirisch
zeigt \citet{zotero_delvecchio_001} für Kalshi-Märkte, dass informierter Handel
real und über eine Abnormal-Trade-Size-Kennzahl detektierbar ist und sich gegen
die Marktauflösung verstärkt.

Für die gebotene Vorsicht sprechen mehrere Befunde. \citet{zotero_poly_001} zeigen
für die Oktober-Grosskonto-Episode 2024, dass Kapital gleichzeitig auf beide
Marktseiten floss --- ein Muster, das mit heterogenen Erwartungen und nicht mit
einseitiger Manipulation vereinbar ist. \citet{zotero_poly_007} ergänzt
Verhaltensverzerrungen, etwa die Ueberschätzung kleiner Wahrscheinlichkeiten,
sowie eine im Vergleich zu traditionellen Instrumenten höhere Volatilität. In der
Praxis setzt Polymarket inzwischen eine On-Chain-Ueberwachung ein, die unter
anderem ungewöhnliche Positionsgrössen vor Ergebnissen flaggt
\citep{news_polymarket_integrity_001}, gestützt auf forensische Heuristiken wie
Funding-Quellen und Wallet-Clustering \citep{news_chainalysis_forensics_001}.
Diese Arbeiten liefern gemeinsam eine Signatur informierten Handels, ohne dass
einzelne Adressen damit als Insider bewiesen wären.

\section{Forschungslücke}
Die Literatur liefert Theorie zur Aggregation und Effizienz, etablierte Methodik
(Brier, Diebold-Mariano, Ereignisstudien, Granger) und gemischte Evidenz zu
Polymarket. Diese Bausteine werden jedoch meist getrennt behandelt: Prognosequalität,
Ereignisreaktion und On-Chain-Timing erscheinen in unterschiedlichen Arbeiten und
selten für denselben Marktkontext. Was fehlt, ist eine reproduzierbare,
deterministisch geprüfte Zusammenführung dieser drei Perspektiven für einen
gemeinsamen Fall. Genau diese Lücke schliesst die vorliegende Arbeit für die
US-Wahl 2024, indem sie H1, H2 und H3 auf derselben Datenbasis und mit
nachvollziehbaren, vorab festgelegten Verfahren untersucht.
